% =============================================================================
%  CHAPTER 5 — EVALUATION
%  SOURCE:   docs/EVALUATION.md + docs/CROSS_AGENT_EVALUATION.md + results/*.json
%  RUBRIC:   Evaluation & Reflection (20%) + heavy Execution (50%) overlap
%  BUDGET:   ~10-14 pages — this chapter earns the most marks per page.
%  TOP BAND: fair baselines, measured (not predicted) numbers, explicit threats
%            to validity, and self-critical interpretation.
%
%  >>> THE #1 RISK (see memory project-evaluation-gap):
%      - Do NOT present the 4 petclinic "PASS (expected)" rows as results.
%      - Do NOT headline scripted-regex baseline vs predicted-structured as the
%        main correctness claim — it is not apples-to-apples.
%      - RUN the structured Claude agent AND the LLM text-edit agent on the
%        9-task petclinic suite, then report THOSE numbers.
%
%  >>> WHAT IS DRAFTED HERE (AI first draft, grounded in your docs — verify every
%      number against results/*.json, then edit into your own voice):
%        - Research Questions, Experimental Setup, Validation Methodology
%        - Sample-project results table (MEASURED: results/structured-4.json vs
%          results/text-edit-1.json)
%        - spring-petclinic scripted baseline (MEASURED 6/9) + root-cause analysis
%        - Threats to Validity (descriptive)
%      WHAT IS YOURS TO WRITE (marked \writeme — the judgement/interpretation):
%        - the Discussion's interpretive claims
%        - the LLM-vs-LLM petclinic numbers slot (run finishing in background)
% =============================================================================
\chapter{Evaluation}
\label{ch:eval}

% --- AI FIRST DRAFT from docs/EVALUATION.md. Edit into your own voice. ---
This chapter evaluates the central claim of the project: that performing
refactorings through the IDE's structured engine, rather than by generating
modified source text, yields transformations that are correct with respect to
the program's semantic model, and does so more efficiently than a text-editing
agent. \Cref{sec:eval-rqs} restates the research questions and maps each to an
experiment. \Cref{sec:eval-setup,sec:eval-validation} describe the benchmark
projects, agents, and the multi-layer validation pipeline that produces the
objective correctness signal. \Cref{sec:eval-sample,sec:eval-petclinic} report
measured results on a small purpose-built project and on the real-world
\texttt{spring-petclinic} application respectively, and \Cref{sec:eval-rootcause}
analyses \emph{why} the text-editing baseline fails the cases it fails.
\Cref{sec:eval-threats} states the threats to validity, and \Cref{sec:eval-discussion}
interprets what the results do and do not establish.

\section{Research Questions}
\label{sec:eval-rqs}
% --- AI FIRST DRAFT. These mirror the RQs you frame in the Introduction; keep
%     the wording consistent with sec:intro-rqs once you write that section. ---
The evaluation is organised around three research questions; a fourth,
concerning developer-perceived effort, is a controlled human study deferred to
future work (\Cref{ch:conclusion}).

\begin{description}
  \item[RQ1 (Correctness).] Does refactoring through the structured IDE engine
    produce semantically correct, compile-safe transformations, and where does a
    text-substitution agent operating on the same tasks fail? \emph{Answered by
    the compile-and-content validation across both benchmark projects
    (\Cref{sec:eval-sample,sec:eval-petclinic}) and the root-cause analysis
    (\Cref{sec:eval-rootcause}).}
  \item[RQ2 (Autonomy).] Can a function-calling LLM, given only the tool API and
    injected project context, complete multi-step refactoring tasks on an
    unfamiliar codebase without manual intervention? \emph{Answered by the
    end-to-end agent runs in which the model is told only the task description and
    must locate symbols and select operations itself.}
  \item[RQ3 (Efficiency).] At equal correctness, does the structured action space
    reduce the interaction cost --- API round-trips and tool invocations --- of
    completing a task relative to raw file editing? \emph{Answered by the turn and
    tool-call measurements in \Cref{sec:eval-sample}.}
\end{description}

\section{Experimental Setup}
\label{sec:eval-setup}
% --- AI FIRST DRAFT from docs/EVALUATION.md §2-5 + CROSS_AGENT_EVALUATION.md. ---

\subsection{Benchmark Projects}
Three Java projects are used. The first, \emph{sample-java-project}, is a Maven
project of ten classes (Java~17) that was constructed specifically to exercise
cross-file refactoring: several tasks rename or move a symbol that is referenced
from other files, so that a tool which updates only the declaration --- and not
its remote references --- produces code that no longer compiles. Because the
project is hand-crafted to expose the structural advantages of AST-aware
refactoring, it functions as a controlled instrument rather than a representative
sample; this is revisited in \Cref{sec:eval-threats}.

The second, \texttt{spring-petclinic}, is the canonical Spring Boot sample
application (cloned at tag \texttt{3.3.0}), comprising roughly thirty classes
organised into the layered packages typical of a real Spring project. It is used
to test whether the agent can navigate an unfamiliar, idiomatic codebase it was
not designed around, and to surface structural hazards --- name collisions across
classes, implicit same-package references, and method overloading --- that a
small artificial project does not naturally contain.

The third, \texttt{commons-lang} (Apache Commons Lang, cloned at tag
\texttt{rel/commons-lang-3.14.0}), is a widely-used utility library of 246
source files --- roughly an order of magnitude larger than the others and from a
different domain. It is included specifically to test whether the structural
findings hold \emph{at scale} on an independent, real-world codebase, addressing
the external-validity concern that the other two projects raise
(\Cref{sec:eval-threats}).

\subsection{Task Suites}
Each project has a corresponding task suite. \texttt{tasks.json} defines eight
tasks on the sample project spanning seven operation types (field and method
rename, move class, safe-delete, inline method, class creation, add method, and
change signature); five of the eight involve at least one cross-file reference.
\texttt{tasks\_petclinic.json} defines nine tasks on \texttt{spring-petclinic}
across the same operation families, including three tasks
(\Cref{sec:eval-rootcause}) specifically chosen to require type- and
reference-level information that text matching cannot recover. Every task carries
a machine-checkable success specification (\Cref{sec:eval-validation}).

\paragraph{Operation coverage and representativeness.} The benchmarked operations
deliberately span both ends of what coding agents do in practice. The large-scale
empirical study of \citet{agenticrefactoring2025} finds that agent refactorings are
dominated by localized renames and type changes (rename-parameter 10.4\%,
rename-variable 8.5\%, change-variable-type 11.8\%), so the rename-heavy core of the
suite mirrors the most common agent behaviour. Conversely, the structural operations
the suite also includes --- move class, change signature, and
overload-disambiguated rename --- are precisely those the same study reports agents
\emph{avoid} relative to human developers, i.e.\ the higher-level changes where
free-text editing is riskiest. The suite therefore exercises an agent both where it
is already active and where a structurally-safe execution layer could expand what it
can attempt without risking incorrect edits.

\subsection{Agents and Modes}
\label{sec:eval-modes}
The primary comparison is between two \emph{edit primitives} --- the structured
IDE engine versus raw text editing --- not between two language models. To
separate the contribution of the edit primitive from that of the driving model's
capability, the primitives are exercised in a controlled $2\times2$ design
(\Cref{tab:modes}): each is driven by (i)~\emph{no LLM} (a scripted ``direct''
runner), (ii)~a \emph{weak, freely-available} model, and (iii)~a \emph{capable
commercial} agent. Holding the model fixed and varying only the primitive
isolates the action space; holding the primitive fixed and varying the model
isolates capability.

\begin{table}[t]
  \centering
  \caption{Evaluation arms. The \emph{direct} modes isolate the edit primitive
  with no LLM; the \emph{agent} modes add an LLM that receives only the task
  description and must locate symbols and choose operations itself. The capable
  agent is the system's intended deployment (a commercial Claude Code agent
  driving the tools); the weak models establish how each action space behaves
  under limited tool-use capability.}
  \label{tab:modes}
  \begin{tabular}{@{}lll@{}}
    \toprule
    Arm & Edit primitive & Driver (model) \\
    \midrule
    Structured direct          & IntelliJ PSI engine  & none (scripted) \\
    Scripted text-edit         & regex/string rewrite & none (scripted) \\
    Structured agent (capable) & IntelliJ PSI engine  & Claude Opus~4.8 \\
    Structured agent (weak)    & IntelliJ PSI engine  & qwen2.5:7b (local) \\
    Text-edit agent (capable)  & raw file I/O         & Claude Opus~4.8 \\
    Text-edit agent (weak)     & raw file I/O         & llama-4-scout (Groq) \\
    \bottomrule
  \end{tabular}
\end{table}

The \emph{direct} modes execute the predefined operation for each task with no
model in the loop: the structured-direct runner calls the plugin's tool API, and
the scripted text-edit runner applies a word-boundary regex substitution. They
isolate the edit primitive itself --- a lower bound for the text side and a
ceiling check for the structured side --- independently of any model's planning.
The scripted baseline is deliberately weaker than a modern LLM agent; it is
included as a transparent lower bound that exposes exactly what raw name-matching
can and cannot do, not as the headline comparator.

The \emph{agent} modes wrap each primitive in an LLM tool-use loop. The agent
receives only the natural-language task description and the injected project
context (\Cref{ch:impl}), and must itself find the target symbol and issue the
appropriate calls; it is never told file paths or text offsets. The runners are
model-agnostic --- a \texttt{--provider} switch drives Anthropic,
OpenAI-compatible, Groq, and local Ollama endpoints --- which both supports the
model-agnostic architectural claim (\Cref{sec:design-mcp}) and lets the
evaluation run on freely available models.

\paragraph{Capable agent (intended deployment).} The capable arm is
\textbf{Claude Opus~4.8} (model identifier \texttt{claude-opus-4-8}). Because no
paid API budget was available, it was exercised in the form the system is
designed for: a commercial \emph{Claude Code} agent driving the tools, realised
as fresh, \emph{blind} subagents --- one per task, given only the task text and
the live tool schema (\texttt{GET /tools/schema}), with no expected answers and
no knowledge of the success criteria. The structured arm makes all edits through
the tools; the text-edit arm edits files directly with no structured tool
available. Both may run \texttt{mvn compile} and iterate to self-correct,
mirroring a real developer's Claude Code session. Each task is independently
scored by the same objective oracle (\Cref{sec:eval-validation}) and the project
is reset between tasks. This use of Claude throughout --- both to orchestrate the
runs and as the agent under test --- is disclosed in the GenAI declaration; it
demonstrates the intended deployment and the action-space ceiling, and is not an
independent third-party benchmark.

\paragraph{Weak agents (capability floor).} The weak structured arm is
\textbf{qwen2.5:7b} served locally through Ollama (no rate limit); the weak
text-edit arm is Groq-hosted \texttt{meta-llama/llama-4-scout-17b-16e-instruct}.
Both use a 20-turn per-task budget at temperature~0. They establish how each
action space degrades under limited tool-use capability --- which, as
\Cref{sec:eval-discussion} shows, differs sharply between the two primitives.

\paragraph{Reproducibility.} The environment was IntelliJ IDEA Community~2025.1.4
(platform build~251), plugin~v0.2.0, launched via Gradle~8.14.3
\texttt{runIde}; Apache Maven~3.8.5 with Oracle JDK~24.0.2 served as the
compile/validation oracle; Python~3.10.5; Ollama~0.30.2. The headless flag
\texttt{-Dide.performance.skip.refactoring.dialogs=true} (and
\texttt{-Didea.trust.all.projects=true} for unattended runs) makes the
refactoring engine non-interactive. Exact commands, model identifiers, and the
per-arm harness are recorded in the project repository
(\texttt{docs/EXPERIMENTAL\_SETUP.md}) so every figure can be reproduced.

\section{Validation Methodology}
\label{sec:eval-validation}
% --- AI FIRST DRAFT from docs/EVALUATION.md §4. ---
Correctness is judged not by inspecting the model's output text but by checking
the state of the project after the operation, through a pipeline of up to four
layers. The key methodological commitment is that \textbf{compilation is the
primary objective correctness signal}: a transformation that misses a cross-file
reference, breaks an import, or renames the wrong symbol manifests as a
\texttt{mvn compile} failure, independently of how plausible the edited text
looks.

\begin{enumerate}
  \item \textbf{Tool-call layer.} Every mutating operation returned success and
    all operations the task requires were actually invoked.
  \item \textbf{Compile layer.} \texttt{mvn compile} succeeds on the project after
    the refactoring. This is the primary signal: it is where an agent that fails
    to update remote references is caught.
  \item \textbf{Disk-content layer.} Task-specific checks on the files on disk:
    that an expected file or symbol exists, or --- for renames and moves --- that
    the \emph{old} symbol name is absent from \emph{all} source files. Content
    checks use a word-boundary regex (\texttt{\textbackslash b<symbol>\textbackslash b})
    so that, for example, \texttt{normalize} does not spuriously match
    \texttt{normalizeInput}. A \emph{surviving-symbol} variant additionally asserts
    that names which \emph{should} have been left untouched are still present,
    which is what catches an over-eager rename (\Cref{sec:eval-rootcause}).
  \item \textbf{RefactoringMiner layer (optional).} For tasks where it applies,
    RefactoringMiner~\citep{refactoringminer2025} is run over the before/after
    commits to confirm that the intended refactoring type (e.g.\ ``Rename Method'',
    ``Move Class'') is detected in the diff.
\end{enumerate}

Both runners enforce \textbf{task isolation}: the project's baseline state is
committed before each task and restored by \texttt{git reset} afterwards, so tasks
cannot contaminate one another. This reset pattern interacts subtly with
IntelliJ's in-memory file model and was the source of one of the engineering
challenges discussed in \Cref{sec:impl-vfs}; the structured runner therefore
forces a virtual-file-system refresh and waits briefly after each reset before
the next task begins.

\section{Results: Sample Project (Measured)}
\label{sec:eval-sample}
% --- AI FIRST DRAFT. Numbers MEASURED and verified against:
%       results/structured-4.json   (structured agent, 8 tasks)
%       results/text-edit-1.json    (text-edit agent, 8 tasks)
%     Recompute before submission if you re-run. ---
On the eight-task sample suite, both agents complete all tasks: the project
compiles and every content check passes in both the structured and the text-edit
modes (\Cref{tab:sample-correctness}). \textbf{On a project of this size,
correctness alone does not separate the two approaches.} A capable LLM driving raw
file edits can compensate for the absence of a reference index by reading every
file and enumerating the call sites by hand --- and on ten classes it has the
context budget to do so.

\begin{table}[t]
  \centering
  \caption{Sample project (8 tasks): correctness and efficiency. Correctness is
  identical; the separation is in interaction cost. Source:
  \texttt{results/structured-4.json}, \texttt{results/text-edit-1.json}.}
  \label{tab:sample-correctness}
  \begin{tabular}{@{}lcc@{}}
    \toprule
    Metric & Structured agent & Text-edit agent \\
    \midrule
    Tasks passed (compile + content) & \textbf{8/8} & \textbf{8/8} \\
    Mean turns (API round-trips) per task     & 3.4 & 4.2 \\
    Mean tool calls per task                  & 3.8 & 9.8 \\
    Median tool calls per task                & 3   & 9.5 \\
    \bottomrule
  \end{tabular}
\end{table}

The separation is instead in \emph{efficiency}. The structured agent completes
each task in fewer API round-trips (mean 3.4 vs 4.2) and substantially fewer tool
calls (mean 3.8 vs 9.8). The mean understates the typical gap, because it is
inflated by a single structured outlier: \texttt{inline-001}, which inlines a
method across five call sites, took 13 structured tool calls. Excluding that task,
the structured agent never exceeds three calls; the \emph{median} of three calls
versus 9.5 is the more representative figure. The mechanism is direct: a
structured rename is one call that updates the declaration and every reference
atomically, whereas the text-edit agent must list files, read several of them,
locate each occurrence, and issue a separate edit per site --- roughly a
$2.6\times$ increase in tool invocations for the same outcome.

For RQ3, this is the trade-off in the structured engine's favour: at equal correctness,
performing the change in a single tool call --- rather than listing files, reading
several of them, and issuing one edit per site --- substantially cuts the agent's
interaction cost. Operationally this matters, because frontier models of the kind behind
Claude Code are metered by tokens and round-trip latency, so fewer calls mean lower cost
and faster completion; and because each discrete operation is a chance to emit a
malformed call or miss a site, fewer operations also mean fewer opportunities to err.

\section{Results: spring-petclinic}
\label{sec:eval-petclinic}
% =============================================================================
%  >>> ACTION REQUIRED before this section is final:
%      1. LLM text-edit petclinic run is executing now ->
%         results/text-edit-petclinic-gptoss.json. Fill tab:petclinic with its
%         per-task PASS/FAIL once it lands.
%      2. Structured-LLM (and remaining structured-direct) petclinic tasks need
%         the sandbox IntelliJ. Until run, DO NOT report them as measured —
%         the table below marks them "not yet run".
% =============================================================================
The sample project shows that a small codebase does not discriminate the two
approaches on correctness. The purpose of \texttt{spring-petclinic} is to test
the correctness claim where it should bite: on an unfamiliar, real-world project
containing structural hazards a ten-class project does not.

\subsection{Scripted text-edit lower bound (measured)}
The scripted text-edit runner --- the no-LLM lower bound --- passes
\textbf{6 of 9} tasks (\Cref{tab:petclinic}). The three failures are not random:
each is a task that requires type or reference information that a word-boundary
text substitution does not have. They are analysed in \Cref{sec:eval-rootcause}.
This establishes, concretely and reproducibly, that raw name-matching is
\emph{insufficient} for three common refactorings on a real codebase --- before
any LLM is involved.

\begin{table}[t]
  \centering
  \caption{spring-petclinic (9 tasks), all columns measured. Scripted text-edit is the
  no-LLM lower bound (\texttt{results/text-edit-petclinic-direct-7.json}); LLM text-edit is
  the \texttt{llama-4-scout} agent (\texttt{results/text-edit-petclinic-llama4scout.json});
  structured executes the plugin's IntelliJ PSI refactoring operations directly, with no
  LLM (\texttt{results/structured-petclinic-direct-FINAL.json}). The structured column is
  the tool-layer ceiling; an LLM-driven structured agent is left to future work
  (\Cref{ch:conclusion}).}
  \label{tab:petclinic}
  \begin{tabular}{@{}llccc@{}}
    \toprule
    Task & Operation & Scr.\ text-edit & LLM text-edit & Structured \\
     & & {\footnotesize(no LLM)} & {\footnotesize(LLM planner)} & {\footnotesize(direct, no LLM)} \\
    \midrule
    \texttt{pc-rename-001}        & rename field      & PASS & PASS & PASS \\
    \texttt{pc-add-method-001}    & add method        & PASS & PASS & PASS \\
    \texttt{pc-find-usages-001}   & find usages       & PASS & \textbf{FAIL}$^\dagger$ & PASS \\
    \texttt{pc-read-file-001}     & read file         & PASS & PASS & PASS \\
    \texttt{pc-rename-002}        & rename class+file & PASS & PASS & PASS \\
    \texttt{pc-rename-method-001} & rename method     & \textbf{FAIL} & PASS & PASS \\
    \texttt{pc-move-001}          & move class        & \textbf{FAIL} & \textbf{FAIL} & PASS \\
    \texttt{pc-rename-method-002} & rename overload   & \textbf{FAIL} & \textbf{FAIL} & PASS \\
    \texttt{pc-change-sig-001}    & change sig.       & PASS & PASS & PASS \\
    \midrule
    \multicolumn{2}{@{}l}{\textbf{Total}} & \textbf{6/9} & \textbf{6/9} & \textbf{9/9} \\
    \bottomrule
  \end{tabular}

  \smallskip
  {\footnotesize $^\dagger$\,The text-edit toolset exposes no structural
  find-usages operation, so the agent cannot answer a project-wide reference query
  by construction.}
\end{table}

\subsection{LLM text-edit agent (measured)}
The LLM text-edit agent shown here --- the \emph{weak} free model
\texttt{llama-4-scout}, driven by raw file tools only (the \emph{capable} text-editing
agent is compared separately in \Cref{sec:eval-controlled}) --- also passes
\textbf{6 of 9} tasks (\Cref{tab:petclinic}), but it is instructive
that it fails a \emph{different} set from the scripted baseline. The agent
\emph{succeeds} on \texttt{pc-rename-method-001}, the cross-class \texttt{addVisit}
collision that defeats the scripted substitution: by reading the relevant files it
renames only \texttt{Owner.addVisit} and its caller while leaving
\texttt{Pet.addVisit} untouched, which the surviving-symbol check confirms. A
capable model can therefore reason its way past a name collision that pure string
matching cannot. It likewise renames the class \emph{and} its file correctly in
\texttt{pc-rename-002}.

It nonetheless fails exactly where text editing lacks information the IDE engine
has by construction. On \texttt{pc-move-001} it moves \texttt{PetValidator} to the
new package (the file appears at the new path) but the project no longer compiles,
because it does not synthesise the now-required import for the same-package
reference to \texttt{Pet}. On \texttt{pc-rename-method-002} it renames the target
overload to \texttt{findPet} but again breaks compilation. And it cannot perform
\texttt{pc-find-usages-001} at all: the text-edit toolset exposes no structural
usage query, so a project-wide reference count lies outside its action space.

The decisive observation is the \emph{intersection of failures}. The two tasks
that defeat \emph{both} a na\"ive script \emph{and} a capable LLM editing raw text
--- \texttt{pc-move-001} and \texttt{pc-rename-method-002} --- are precisely those
that require synthesising an implicit import and performing an overload-safe rename
with caller updates. These are not the kind of failure a larger model would
obviously remove; they are failures of the \emph{edit primitive}, which acts on
text rather than on a typed, project-wide reference graph
(\Cref{sec:eval-rootcause}).

\paragraph{A practical constraint worth recording.} Running an LLM text-edit agent
over a real project is token-intensive: each turn resends the growing transcript of
file contents the agent has read. A first attempt using a larger model
(\texttt{gpt-oss-120b}) on a free API tier could not complete the suite at all ---
its per-minute token limit (8{,}000) is exceeded once one or two source files are
in context, so six of nine tasks aborted before any edit was made. The run reported
here used a model whose free-tier per-minute budget (30{,}000 tokens) was large
enough to sustain the loop. The structured agent never reads whole files into the
prompt to act on them (\Cref{sec:eval-sample}), so its per-task token footprint is
far smaller --- an efficiency advantage with a direct cost dimension.

\subsection{Structured operations (measured)}
Executing the same nine tasks through the plugin's IntelliJ refactoring operations
yields \textbf{9 of 9} (\Cref{tab:petclinic}). Decisively, the structured engine
passes every task that defeats one or both text approaches, and the validation
signals confirm \emph{why}:
\begin{itemize}
  \item \texttt{pc-move-001}: \texttt{PetValidator} is moved to the \texttt{system}
    package and the project still \emph{compiles} --- \texttt{MoveClassesOrPackagesProcessor}
    synthesised the import for the same-package reference to \texttt{Pet} that the
    text agents could not (\Cref{sec:eval-rootcause}).
  \item \texttt{pc-rename-method-002}: only the \texttt{(String, boolean)} overload
    becomes \texttt{findPet}; the surviving-symbol check confirms the other
    \texttt{getPet} overloads are untouched --- the type-signature targeting that
    text matching lacks.
  \item \texttt{pc-rename-method-001}: \texttt{Owner.addVisit} becomes
    \texttt{recordVisit} while \texttt{Pet.addVisit} survives, and
    \texttt{pc-rename-002} renames the class \emph{and} its file on disk.
\end{itemize}
Against these baselines the headline is \textbf{structured 9/9 versus 6/9 for both
text approaches}, with the gap concentrated in exactly the three reference-graph /
type-system operations of \Cref{sec:eval-rootcause}. Two qualifications are essential,
and are developed in \Cref{sec:eval-controlled}. First, this column isolates the
\emph{edit primitive}: it executes the operations directly, with no LLM planner, and so
measures whether the structured tools \emph{can} perform each task correctly --- which
they do by construction. That an LLM can itself \emph{select and invoke} these
operations correctly is confirmed by running the structured \emph{agent} (Claude
Opus~4.8) end-to-end: it also scores 9/9 here, and 4/4 and 3/3 on the two larger
projects (\Cref{sec:eval-controlled}) --- so the LLM-driven structured agent is now
\emph{measured}, not deferred. Second, the ``6/9 for text'' figure is for the scripted
baseline and a \emph{weak} free model; a \emph{capable} text-editing agent does
markedly better (7/9), so the durable claim is not raw correctness against any text
editor but the action space's behaviour under limited capability, together with the
structural queries and precision that text editing cannot match.

\paragraph{Engineering the structured runner to be non-interactive (and a threat to
validity).} Driving IntelliJ's refactoring engine unattended is not free: the engine
is built for an interactive user, so multi-file operations raise modal dialogs (a
usage preview, a conflicts dialog, an automatic-rename dialog) that block the UI
thread indefinitely when no one is present to dismiss them. Reaching the 9/9 above
required making the operations genuinely headless --- routing them through the
platform's non-interactive code paths (the \texttt{ide.performance.skip.refactoring.dialogs}
switch and proceeding past recorded conflicts) and ensuring each operation resolves
an in-project element after reconciling the virtual file system with the on-disk
state. This is the same external-change hazard analysed in \Cref{sec:impl-vfs},
encountered again here: the benchmark git-resets the project between tasks, and a
stale virtual file system manifested as refactorings that executed in the PSI but
did not persist, or as ``element is not located inside the project'' precondition
errors. This is on-thesis rather than incidental: the structured numbers were obtained
on a freshly imported, isolated IDE instance, and making the headless path robust ---
suppressing the modal dialogs and reconciling the virtual file system after each
git-reset (\Cref{sec:impl-vfs}) --- was itself an engineering result rather than
something the platform provided for free. It is a fair counter-cost to record against
the approach: the structured engine's correctness guarantee comes with the burden of
driving a UI-oriented IDE unattended.

\subsection{Diff Minimality}
\label{sec:eval-diff}
% --- AI FIRST DRAFT. Numbers measured + instrumented:
%       results/structured-petclinic-direct-DIFF.json  (structured)
%       results/text-edit-scripted-DIFF.json           (scripted text-edit)
%       results/text-edit-petclinic-DIFF.json          (LLM text-edit; only 2 tasks
%         completed before the Groq free-tier daily token cap). ---
Beyond pass/fail, the \emph{size} of the change a refactoring produces matters for
reviewability; small, surgical diffs are easier to audit, and
\citet{agenticrefactoring2025} note that agent refactorings are valued for being
localized. Instrumenting the runners to record the source-line delta of each task
(\Cref{tab:diff}) shows structured execution is consistently the more surgical.
Across the seven mutating petclinic tasks the structured engine changes a median of
8 lines (mean 8.3, 58 in total), against the scripted text-edit baseline's median 12
(mean 21.7, 152 in total --- roughly $2.6\times$ more); the LLM text-edit agent, on
the two tasks it completed before exhausting the free API tier's daily token budget,
likewise rewrote more (17 vs 10 lines for the field rename, 6 vs 3 for add-method).

\begin{table}[t]
  \centering
  \caption{Source-line delta per task (lines changed\,/\,files touched). Structured
  and scripted text-edit are complete; the LLM text-edit agent (\texttt{llama-4-scout})
  completed only two tasks before the free-tier daily token cap (the rest are omitted
  rather than counted as zero). Read-only tasks (find-usages, read-file) change
  nothing and are excluded.}
  \label{tab:diff}
  \begin{tabular}{@{}lccc@{}}
    \toprule
    Task & Structured & Scripted text-edit & LLM text-edit \\
    \midrule
    \texttt{pc-rename-001}        & 10\,/\,1 & 12\,/\,1            & 17\,/\,2 (fail) \\
    \texttt{pc-add-method-001}    & 3\,/\,1  & 1\,/\,1            & 6\,/\,1 \\
    \texttt{pc-rename-002}        & 8\,/\,3  & 37\,/\,1           & --- \\
    \texttt{pc-rename-method-001} & 6\,/\,3  & 10\,/\,3 (fail)    & --- \\
    \texttt{pc-move-001}          & 5\,/\,3  & 64\,/\,1 (fail)    & --- \\
    \texttt{pc-rename-method-002} & 8\,/\,2  & 20\,/\,3 (fail)    & --- \\
    \texttt{pc-change-sig-001}    & 18\,/\,6 & 8\,/\,2           & --- \\
    \midrule
    \textbf{Total} & \textbf{58} & \textbf{152} & --- \\
    \bottomrule
  \end{tabular}
\end{table}

The contrast is sharpest on the structural tasks. On \texttt{pc-move-001} the
scripted approach rewrites 64 lines in a single file and still fails to compile,
whereas the structured move touches 5 lines across 3 files and succeeds; on the
overload rename the structured engine changes 8 lines (correct) against the
baseline's 20 (incorrect). The one task where the structured engine changes
\emph{more} lines is \texttt{pc-change-sig-001} (18 lines across 6 files): correctly
propagating a new parameter to every call site legitimately requires touching every
caller. This is the right reading of minimality --- not the fewest lines in the
absolute, but the minimal \emph{correct} change. Structured edits are confined to a
symbol's declaration and its genuine references; text approaches make incidental,
over-broad substitutions that touch more code and, on the structural tasks, break
it. This also bears on \emph{auditability}: a small, structurally-scoped diff, together
with the explicit tool-call trace the agent produces (\Cref{ch:design}), is far easier
for a developer to review and trust than a large free-text rewrite whose intent must be
reverse-engineered from the changed lines.

\subsection{Root-Cause Analysis of the Text-Edit Failures}
\label{sec:eval-rootcause}
% --- AI FIRST DRAFT from docs/EVALUATION.md §7b. This material is strong and
%     fully grounded in the task definitions; verify file/line references. ---
The three scripted failures share a single underlying cause: text-level editing
lacks the reference graph and the type information needed to act on the right
symbol. Each failure isolates one facet of that deficiency.

\paragraph{Cross-class name collision (\texttt{pc-rename-method-001}).}
The task renames \texttt{Owner.addVisit(Integer, Visit)} to \texttt{recordVisit}.
A different class, \texttt{Pet}, independently declares \texttt{addVisit(Visit)}.
The scripted baseline performs a word-boundary substitution of \texttt{addVisit},
which cannot distinguish the two declarations: it also renames \texttt{Pet.addVisit}
and the \texttt{pet.addVisit(visit)} call in \texttt{VisitController}, violating
the task. IntelliJ's \texttt{RenameProcessor} operates on the specific PSI node
for \texttt{Owner.addVisit}; its reference index knows that
\texttt{pet.addVisit(visit)} resolves to \texttt{Pet.addVisit} and leaves it
untouched. The \emph{surviving-symbol} check detects the difference: after the
structured rename, \texttt{addVisit} still appears as \texttt{Pet}'s declaration
and at its call site, whereas after the text-edit rename it has been erased from
every file.

\paragraph{Implicit same-package dependency (\texttt{pc-move-001}).}
\texttt{PetValidator} and \texttt{Pet} both live in the \texttt{owner} package, so
\texttt{PetValidator} references \texttt{Pet} with no \texttt{import} statement;
likewise \texttt{PetController} uses \texttt{PetValidator} without one. Moving
\texttt{PetValidator} to the \texttt{system} package requires two new imports to be
synthesised. The scripted baseline can only rewrite \emph{explicit}, fully
qualified references, so it leaves both files importless and the project fails to
compile (\texttt{cannot find symbol: class Pet}). IntelliJ's
\texttt{MoveClassesOrPackagesProcessor} resolves references through the PSI graph,
which includes same-package references that have no import line, and inserts the
required imports automatically.

\paragraph{Overload collision (\texttt{pc-rename-method-002}).}
\texttt{Owner} declares three overloads of \texttt{getPet}: \texttt{(String)},
\texttt{(Integer)}, and \texttt{(String, boolean)}. The task renames only the
\texttt{(String, boolean)} overload to \texttt{findPet}. The scripted substitution
matches the name \texttt{getPet} and renames all three; the project still compiles
(the renames are internally consistent), so this failure would be \emph{invisible
to a compile-only check} --- it is caught only by the surviving-symbol assertion
that \texttt{getPet} must still be present. IntelliJ resolves the target by typed
signature (\texttt{Owner\#getPet(String,boolean)}) and renames exactly one
overload. This case is the clearest illustration of why compilation alone is a
necessary but not sufficient correctness signal.

\paragraph{Synthesis.} In each case the structured engine succeeds because it acts
on a typed node in a project-wide reference graph, while the text substitution acts
on a name string. The deficiency is not that the baseline is poorly engineered but
that the information it would need --- which class a call belongs to, which
references resolve without an import, which overload a signature selects --- is
simply not present in the text. This is the core argument of the thesis made
concrete: see \Cref{sec:eval-discussion}.

\section{Results: Apache Commons Lang (Second Real Project)}
\label{sec:eval-commons}
% --- AI FIRST DRAFT. Measured: results/commons-lang-structured-DIFF.json (structured,
%     no-LLM, via the sandbox IDE) vs results/commons-lang-scripted-DIFF.json (scripted
%     text-edit). 4-task suite: benchmarks/tasks_commons_lang.json. ---
To test whether the structural findings generalise beyond \texttt{spring-petclinic},
the comparison was repeated on Apache Commons Lang --- an order of magnitude larger
(246 source files) and from an unrelated domain (a general-purpose utility library).
A focused four-task suite targets the two discriminating operations
(\Cref{tab:commons}): an overload-disambiguated rename (rename only the
three-argument \texttt{StringUtils.replace} overload), a cross-class collision
(rename \texttt{ArrayUtils.contains(int[],int)} while \texttt{StringUtils.contains}
must survive), plus a field rename and a find-usages query as controls.

\begin{table}[t]
  \centering
  \caption{Apache Commons Lang (4 tasks), measured. Structured executes the IntelliJ
  PSI operations; scripted text-edit is the no-LLM regex baseline. ``Lines'' is the
  source-line delta of the change (\Cref{sec:eval-diff}).}
  \label{tab:commons}
  \begin{tabular}{@{}llcc@{}}
    \toprule
    Task & Operation & Structured & Scripted text-edit \\
     & & {\footnotesize(direct, no LLM)} & {\footnotesize(no LLM)} \\
    \midrule
    \texttt{cl-rename-field-001}       & rename field       & PASS (32 lines)  & PASS (480 lines) \\
    \texttt{cl-rename-overload-001}    & overload rename    & PASS (356 lines) & \textbf{FAIL} (compile) \\
    \texttt{cl-rename-cross-class-001} & cross-class rename & PASS (422 lines) & \textbf{FAIL} (compile) \\
    \texttt{cl-find-usages-001}        & find usages        & PASS (213 refs)  & PASS (40 grep hits) \\
    \midrule
    \multicolumn{2}{@{}l}{\textbf{Total}} & \textbf{4/4} & \textbf{2/4} \\
    \bottomrule
  \end{tabular}
\end{table}

The outcome replicates the petclinic finding on a far larger, independent codebase:
the structured engine passes all four tasks, while the scripted text-edit baseline
fails both structural tasks at the compile stage. The mechanism is identical to
\Cref{sec:eval-rootcause} --- a word-level substitution of \texttt{replace} also
rewrites \texttt{java.lang.String\#replace} calls, and \texttt{contains} rewrites JDK
\texttt{List}/\texttt{String} calls, so the project no longer compiles --- but the
scale is greater: \texttt{ArrayUtils.contains(int[],int)} is referenced from 56 files,
each of which the structured rename updates correctly and atomically.

Two observations are specific to scale. First, the find-usages query exposes the gap
between structural and textual reference-finding directly: the structured engine
resolves \textbf{213} genuine references to \texttt{CharRange}, whereas the baseline's
grep returns only 40 textual matches --- the structural index is both more complete
and free of false positives. Second, diff size behaves differently from the small
project. On the field rename the structured edit is surgical (32 lines) while the
regex baseline changes \textbf{480} lines --- it substitutes every lexical occurrence
of \texttt{calendar} across the library and still compiles, a vivid illustration of
how a text approach can pass a naive check while making a sweeping, semantically
unwarranted change. For the ubiquitous \texttt{replace} and \texttt{contains} methods,
by contrast, the structured edit is itself large (356 and 422 lines): correctly
updating every caller of a widely-used method legitimately touches many files. The
robust claim at scale is therefore not ``fewer lines'' but \emph{minimal correct
change} --- the structured engine changes exactly the genuine references and nothing
else, whereas the text approach changes a comparable or larger amount incorrectly.
This second real-world project strengthens external validity: the structural
discriminators that defeat text editing replicate on an independent codebase an order of
magnitude larger, and the controlled comparison (\Cref{sec:eval-controlled}) extends the
evidence to a third, larger project (\texttt{jackson-databind}, 474 files). The central
findings therefore rest on three independent real-world codebases, not one.

\section{Controlled Comparison: Holding the Model Constant}
\label{sec:eval-controlled}
The comparisons above pair the structured engine against a \emph{weaker} text side
--- a no-LLM regex script, or a small free model. That leaves a confound: does the
structured engine win because of the action space, or because the text side had a
weaker planner? The final experiment removes the confound by holding the model
fixed --- the capable agent, Claude Opus~4.8 (\Cref{sec:eval-modes}) --- and varying
only the edit primitive, on both real-world projects.

\begin{table}[t]
  \centering
  \caption{Controlled comparison: the \emph{same} capable agent (Claude Opus~4.8)
  driving each edit primitive, scored by the identical oracle. The last column
  excludes the non-refactoring probes --- a project-wide find-usages query in each
  project, plus one read probe --- that text editing cannot express as a structured
  tool call.}
  \label{tab:controlled}
  \begin{tabular}{@{}lccc@{}}
    \toprule
    Project (tasks) & Structured & Text-edit & Refactoring tasks only \\
    \midrule
    spring-petclinic (9)  & 9/9            & 7/9            & 7/7 vs.\ 7/7 \\
    Commons Lang (4)      & 4/4            & 3/4            & 3/3 vs.\ 3/3 \\
    jackson-databind (3)  & 3/3            & 2/3            & 2/2 vs.\ 2/2 \\
    \midrule
    \textbf{Total}        & \textbf{16/16} & \textbf{12/16} & \textbf{12/12} \\
    \bottomrule
  \end{tabular}
\end{table}

This result is reported in full because it \emph{qualifies the headline}. With a
capable model in the loop, the two action spaces are \textbf{indistinguishable on
behaviour-preserving refactoring correctness}: on every rename, move, and
change-signature task --- including the overload-disambiguated and cross-class
renames that defeat the regex baseline, and at the scale of \texttt{jackson-databind}
(a class referenced from 107 files) --- the agent editing raw text reaches the same
compiling, content-correct result as the agent driving the IDE engine. It does so by
supplying for itself the very capability the structured engine externalises: reading
the affected files, enumerating the call sites, and re-compiling to catch its own
mistakes.

The entire measured gap (16/16 vs.\ 12/16) lies in the tasks where text editing has
\emph{no corresponding primitive}: a project-wide \texttt{find\_usages} query in each
of the three projects --- where the text agent can only grep, returning, on Commons
Lang, 40 lexical hits against the structured index's 213 genuine, type-aware
references --- plus one read probe scored on a specific tool call. \textbf{On the
twelve behaviour-preserving refactoring tasks the two action spaces are tied, 12/12}
--- including \texttt{jackson-databind}'s 551- and 349-reference class renames, which
the text agent completed by hand. Against a capable text editor, then, the structured
engine's advantage on \emph{correctness} narrows to the structural \emph{queries} that
text fundamentally cannot perform.

This sharpens, rather than collapses, the thesis. It locates precisely where the
structured action space earns its value:
\begin{itemize}
  \item \textbf{Driver capability.} The correctness advantage is decisive exactly when
  the driver \emph{cannot} compensate by reading and re-compiling. The no-LLM regex
  baseline fails the structural tasks outright (2/4 on Commons Lang,
  \Cref{tab:commons}); it is the capable model's iterative self-correction --- not the
  text primitive itself --- that closes the gap.
  \item \textbf{Correct by construction vs.\ contingently correct.} The structured
  engine is correct in a single pass (it scored 7/9 here even \emph{without} a feedback
  loop). The text agent's parity was obtained \emph{with} a compile-and-fix loop;
  where no build feedback is available, that contingency is what fails.
  \item \textbf{Structural queries.} Exhaustive, type-aware reference finding has no
  text-level equivalent, and is a primitive real refactoring workflows depend on.
  \item \textbf{Operation count and precision.} The structured agent completes each
  refactoring in a single mutating call that touches only the typed references. A
  capable text agent is, perhaps surprisingly, comparably economical in raw
  \emph{operations} --- it renamed the 551-reference \texttt{ClassUtil} across 107
  files with two bulk substitutions and two compile checks. But a bulk substitution
  matches by \emph{name}: it can also rewrite comments, strings, and unrelated
  identifiers that share the name, and ``it compiled'' does not establish that it did
  not (cf.\ the regex baseline's 480-line over-broad rename, \Cref{tab:commons}). The
  structured edit is precise by construction; the text edit is a heuristic that
  happened to be correct here.
\end{itemize}

\section{Threats to Validity}
\label{sec:eval-threats}
% --- AI FIRST DRAFT. Required for the top band; verify each claim holds for your
%     final runs and add/remove threats honestly. ---
\paragraph{Construct validity.} The sample project was \emph{deliberately designed}
to expose cross-file weaknesses of text editing; it is an instrument for isolating
the phenomenon, not a representative codebase, and its results should be read as
such. The \texttt{spring-petclinic} suite mitigates this by using an external,
real-world project the author did not write, but the nine tasks were still chosen
by the author and are not a random sample of maintenance work.

\paragraph{Internal validity.} The scripted text-edit baseline is an intentional
lower bound and does not, on its own, support a claim about LLM agents; the LLM
text-edit agent (\Cref{sec:eval-petclinic}) is the fair comparator and the basis
for any claim about agents. Correctness is measured by compilation plus content
checks rather than by behavioural (test-suite) equivalence, so a change that
compiles and passes the content checks but alters runtime behaviour would not be
caught; the operations evaluated are, however, all behaviour-preserving by
construction in the structured engine.

\paragraph{External validity.} Four projects and one language (Java; Kotlin
symbol-creation is implemented but not benchmarked) are evaluated, over a bounded set
of operation types. Three are real-world codebases the author did not write
(\texttt{spring-petclinic}; Apache Commons Lang at 246 source files; and
\texttt{jackson-databind} at 474), and the controlled capable-agent comparison
(\Cref{sec:eval-controlled}) was run on all three --- so the central findings are
\emph{measured} on independent projects at scale, including class renames touching
107 files, not merely hypothesised. Two caveats remain. The tasks were chosen by the
author rather than sampled randomly from real maintenance history. And correctness is
judged by compilation plus content checks, not behavioural equivalence; in particular,
a capable text agent's bulk substitution can be \emph{over-broad} (rewriting comments
or strings) yet still compile, so the structured engine's precision advantage
(\Cref{sec:eval-controlled}) is only partly captured by a pass/fail compile oracle and
is evidenced separately by diff size (\Cref{sec:eval-diff}).

\paragraph{Model validity.} Agent results may be specific to the models used. The
capable agent is Claude Opus~4.8; the weak arms are qwen2.5:7b (structured) and
\texttt{llama-4-scout} (text-edit). The qualitative pattern was consistent across the
three projects --- the capable agent reaches parity on refactorings via either
primitive, while the action space is decisive for the weak/scripted drivers --- but it
rests on a single capable model.
That said, the architectural claim does not depend on Claude specifically: the tool API
is model-agnostic (\Cref{sec:design-mcp}), so any function-calling model can drive it,
and re-running the capable arm on a second frontier model (for example a GPT-class model
through the OpenAI-compatible provider) is the natural way to confirm the observed
pattern is not specific to one model (\Cref{ch:conclusion}).

\section{Discussion}
\label{sec:eval-discussion}
% --- This section is YOUR interpretation. A scaffold of the factual frame is
%     drafted; the claims and their hedging are yours to own. ---
The measured results support a \emph{refined} version of the original thesis, and
it is worth stating precisely what they do and do not show.

Against the baselines an automated system would otherwise fall back on --- a scripted
text editor, or a small free model --- the structured engine wins clearly on
correctness: 9 of 9 versus 6 of 9 on spring-petclinic, and the same pattern on Apache
Commons Lang (4/4 versus 2/4), with the gap sitting exactly on the reference-graph and
type-dependent tasks. For drivers of that kind --- which is the regime of cheap, local,
or rate-limited automation --- the action space is decisive.

The controlled comparison then adds the honest qualifier. Given a \emph{capable} agent
--- Claude Opus~4.8, with a compile-and-fix loop --- raw text editing reaches the same
compiling, content-correct result on every behaviour-preserving refactoring, tied 12 of
12 across the three projects, even at the scale of jackson-databind, where it renamed a
551-reference class by hand. I therefore do not claim raw-correctness superiority for
the structured engine over a frontier model: the data does not support it, and stating
so plainly is part of the contribution.

What the data \emph{does} support is narrower but solid. First, a project-wide,
type-aware structural query such as exhaustive \texttt{find\_usages} has no text-level
equivalent --- it is the one task the text agent fails in every project, including the
capable run at scale (2 of 3 on jackson-databind, the single miss being precisely that
query). Second, the structured edit is correct \emph{by construction}: it touches only
typed references, never comments or strings, and needs no build-feedback loop, whereas
the capable text agent's parity was \emph{contingent} on having one. Third, it reaches
the result in far fewer, semantically-scoped operations. Crucially, these advantages are
largest exactly where the driving agent is weakest --- scripted, small, or unverified
--- so the value of the structured action space rises as the driver's capability, or its
access to feedback, falls.

One dimension this evaluation leaves open is cost. The efficiency advantage is argued
here from operation counts and token footprint rather than measured directly; a natural
next step is to quantify the tokens consumed and the wall-clock cost of each approach
across several models, which would establish how large the structured engine's cost
advantage actually is and whether it persists as models improve.
